\section{Tool Comparison and Selection}

The external tools used in this thesis were selected based on the requirements derived in the previous section. 
The comparison is therefore not intended as a comprehensive survey of all available modeling and simulation environments.
Instead, it focuses on those tools that were relevant for the required workflow and that were considered realistic candidates for integration into \emph{SysSimX}.
For each modeling domain, the selected tool is justified with respect to the required capabilities, the integration effort, and the overall scope of the thesis.

\subsection{Equation-Based Modeling}

For the equation-based part of the workflow, the main distinction is between signal-based and equation-based modeling paradigms. 
Signal-based modeling describes systems as directed chains of inputs and outputs.
It is well suited for control logic, signal processing, and explicitly causal subsystem descriptions.
Equation-based modeling, in contrast, describes the relations between physical quantities declaratively through equations.
The computational direction is not fixed in advance, but is determined during model translation.
This makes equation-based modeling particularly suitable for multi-domain physical systems with constraints and conservation laws.

Within the considered tools, Simulink represents the signal-based approach, whereas Modelica and Simscape support acausal physical modeling.
Simulink is based on directed block diagrams and is therefore natural for controllers, signal-flow architectures, and sampled logic ~\cite{noauthor_simulink_nodate}.
Simscape extends this environment by adding acausal physical components and physical network semantics. However, Simscape remains embedded in the proprietary MATLAB/Simulink ecosystem and relies on explicit interface blocks when physical networks interact with signal-based control models ~\cite{noauthor_simscape_nodate}.

Modelica follows a more general equation-based approach.
It provides a dedicated language for object-oriented, multi-domain, and acausal modeling of physical systems. 
In contrast to signal-based block modeling, subsystem behavior is formulated directly in terms of equations and connector relations ~\cite{Modelica_Fritzson}.
This is well aligned with the requirements of the present thesis, since the workflow must support physically meaningful subsystem models, hierarchical composition, and later integration into a heterogeneous co-simulation framework.

For the present work, Modelica was selected as the modeling formalism for equation-based subsystem models, and OpenModelica was selected as the corresponding open-source toolchain.
The decisive reasons were the explicit support for equation-based and acausal modeling, the suitability for multi-domain physical systems, the availability of an open-source compiler, and the possibility to integrate the resulting models into the overall workflow through standard co-simulation interfaces.
Simulink was considered relevant as a representative signal-based environment, and Simscape as an acausal physical extension within that environment.
However, the thesis required an open and scriptable workflow centered on equation-based physical modeling rather than a proprietary block-diagram environment. For this reason, Modelica/OpenModelica was the more suitable choice.

\begin{table}[htbp]
\centering
\small
\caption{Comparison of equation-based modeling candidates.}
\label{tab:eq_tool_comparison}
\begin{tabular}{p{2.2cm} p{2.5cm} p{5.0cm} p{4.0cm}}
\toprule
\textbf{Tool} & \textbf{Paradigm} & \textbf{Relevant Strengths} & \textbf{Main Limitation} \\
\toprule
Simulink
& Signal-based, causal
& \begin{itemize}[nosep,leftmargin=*,topsep=0pt]
    \item Control logic and signal processing
    \item Hierarchical block models
    \item Wide industry adoption
  \end{itemize}
& Not centered on equation-based physical modeling \\
\toprule
Simscape
& Acausal physical (within Simulink)
& \begin{itemize}[nosep,leftmargin=*,topsep=0pt]
    \item Physical network modeling
    \item Combination with Simulink control
  \end{itemize}
& Proprietary MATLAB/Simulink ecosystem \\
\toprule
\textbf{Modelica}
& Equation-based, acausal
& \begin{itemize}[nosep,leftmargin=*,topsep=0pt]
    \item Multi-domain physical modeling
    \item Declarative equation formulation
    \item Open-source toolchain (OpenModelica)
    \item FMI export for co-simulation
  \end{itemize}
& Requires external toolchain for translation and execution \\
\bottomrule
\end{tabular}
\end{table}


% \begin{table}[h]
% \centering
% \caption{Relevant comparison of equation-based modeling candidates.}
% \begin{tabular}{p{2.3cm} p{2.8cm} p{5.2cm} p{4.2cm}}
% \toprule
% \textbf{Tool} & \textbf{Paradigm} & \textbf{Relevant strengths} & \textbf{Main limitation for this thesis} \\
% \midrule
% Simulink & Signal-based, causal & Strong for control logic, signal processing, and hierarchical block models & Not centered on equation-based physical modeling \\
% Simscape & Acausal physical modeling within Simulink & Physical network modeling, combination with Simulink control models & Embedded in proprietary MATLAB/Simulink environment \\
% Modelica & Equation-based, acausal & Multi-domain physical modeling, declarative equations, open-source toolchain, suitable for co-simulation export & Requires external toolchain for translation and execution \\
% \bottomrule
% \end{tabular}
% \end{table}

\subsection{Musculoskeletal Modeling}

For the musculoskeletal part of the workflow, the selected tool had to support articulated biomechanical models with joints, constraints, actuators, and forward dynamic simulation.
In addition, the tool had to be accessible from Python and suitable for integration into a custom co-simulation framework.
Since the thesis focuses on heterogeneous system integration rather than on benchmarking musculoskeletal modeling environments, the goal of this subsection is not to provide a comprehensive survey of all available tools, but to justify the selected tool with respect to the required capabilities.

OpenSim was selected for this purpose.
It is an open-source software platform for modeling, simulating, and analyzing musculoskeletal systems and dynamic movement.
OpenSim provides forward dynamic simulation, a component-based biomechanical model structure, and programmatic access through its application programming interface ~\cite{OpenSim_Delp, seth_opensim_2018}.
These properties make it suitable for representing biomechanical subsystems within the heterogeneous workflow considered in this thesis.

From the perspective of framework integration, OpenSim offers three decisive advantages.
First, it provides an established modeling environment for rigid-body and musculoskeletal dynamics, including joints, constraints, actuators, and contact-related elements.
Second, it supports forward simulation of biomechanical models, which is required for coupling the subsystem to controllers and external components. Third, its Python interface enables direct integration into the \emph{SysSimX} architecture without requiring an intermediate proprietary execution environment.

Other musculoskeletal modeling tools exist, including both research-oriented and commercial environments. However, a systematic comparison of such tools was outside the scope of this thesis. The selection of OpenSim was therefore based on its suitability for the required workflow, its open-source availability, and its practical integrability into the Python-based co-simulation framework.

\subsection{Finite Element Modeling}

For the finite element modeling part of the workflow, the selected tool had to support spatially discretized models, programmable definition of material laws and boundary conditions, and a level of flexibility that allows integration into a custom research-oriented co-simulation framework.
In addition, the tool had to be accessible from Python, since the overall architecture of \emph{SysSimX} is implemented in Python and relies on backend adapters for coupling heterogeneous subsystem models.
As in the musculoskeletal case, the purpose of this subsection is not to provide a comprehensive benchmark of finite element tools, but to justify the selected tool with respect to the requirements of the thesis.

Netgen/NGSolve was selected for this purpose. It is an open-source finite element environment that combines mesh generation with a programmable finite element library ~\cite{schoberl_c11_nodate, schoberl_netgen_1997}. In contrast to solver environments that are primarily organized around graphical pre-processing and fixed analysis workflows, NGSolve provides direct programmatic access to the mathematical problem formulation, including finite element spaces, weak forms, loads, materials, boundary conditions, grid functions, and solver configuration.
This makes it well suited for integrating finite element models, such as  continuum-mechanical subsystem models, into a custom co-simulation framework.

From the perspective of framework integration, the decisive advantage of Netgen/NGSolve is its Python interface. It allows the finite element model to be assembled, parameterized, advanced, and queried directly from the same programming environment in which the co-simulation framework is implemented.
This is particularly useful for the present work, where the finite element subsystem must be wrapped as a co-simulation component and coupled to equation-based and musculoskeletal subsystem models through a common interface.

The selected tool also reflects the scope of the thesis.
The goal was to demonstrate the integration of a finite element subsystem into a heterogeneous hybrid simulation workflow, not to perform an extensive comparative study of finite element platforms.
Other open-source tools such as CalculiX~\cite{noauthor_calculix_nodate}, FEBio~\cite{maas_febio_2012}, or openCFS~\cite{schoder_opencfs_2025} were considered at a high level.
However, for the present work, their integration paths were assessed as less direct for a Python-centered prototype workflow, for example due to less mature programmatic interfaces or a stronger focus on file-based pre-processing and execution.
Netgen/NGSolve was therefore chosen because it combines strong finite element capabilities with direct scriptability and practical suitability for framework integration.

A limitation of this choice is that NGSolve requires a comparatively strong mathematical understanding of finite element modeling, since many aspects of the model must be defined directly in code.
This limitation was accepted because the thesis focuses on framework integration and controllable programmatic access rather than on interactive finite element model authoring.

\subsection{Summary}

In summary, the selected toolchain covers the three modeling domains required in this thesis: 

\begin{itemize}
    \item Modelica/OpenModelica for equation-based physical modeling,
    \item OpenSim for musculoskeletal biomechanics, and
    \item Netgen/NGSolve for continuum-mechanical finite element modeling.
\end{itemize}

Each tool was chosen because it satisfies the relevant domain-specific requirements while remaining sufficiently accessible for integration into a Python-based workflow.
At the same time, no single existing environment covered the full combination of heterogeneous model integration, hybrid co-simulation, and unified orchestration required for this work.
This gap motivated the development of the \emph{SysSimX} framework presented in the following chapters.